\begin{lemma}
Let $n\ge1$ be an integer and let $y\colon\mathbb{R}\to\mathbb{C}$ be
$C^\infty$ on $(0,1)$ and satisfy
\[
  r\bigl(r\,y'\bigr)'(r)+2n\,r\,y'(r)+n^2\,y(r)=0
  \qquad\text{for every }r\in(0,1).
\]
Then there are $A,B\in\mathbb{C}$ with
\[
  y(r)=\bigl(A+B\log r\bigr)r^{-n}\qquad\text{for every }r\in(0,1).
\]

This is the equation (35) of the paper at the single pair $(k,p)=(0,2)$, where
$pn=2n$ and $(n^2-k^2)(p-1)=n^2$.  It is exactly the case in which the
indicial polynomial $\lambda^2+pn\lambda+(n^2-k^2)(p-1)$ degenerates to
$(\lambda+n)^2$, so that the two roots coincide at $-n$ and the paper's
assertion at line 313 that they are distinct fails; the solution
$(\log r)\,r^{-n}$ obtained here for $A=0$, $B=1$ is not a combination of
powers $r^{\lambda}$ at all.
\end{lemma}
\begin{proof}
First note the degeneration of the indicial polynomial recorded in the
statement: at $k=0$ and $p=2$,
\[
  \lambda^2+pn\lambda+(n^2-k^2)(p-1)=\lambda^2+2n\lambda+n^2=(\lambda+n)^2 ,
\]
so $-n$ is a double root.  This is the algebraic reason a logarithm appears
below, and the reason the two-power conclusion is unavailable.

\emph{Step 1: the substitution $r=e^{t}$.}  Put $J=(-\infty,0)$ and
\[
  Y(t)=y\bigl(e^{t}\bigr)\qquad(t\in J).
\]
For $t<0$ one has $e^{t}\in(0,1)$, so $Y$ is defined on $J$, and $Y$ is
$C^\infty$ on $J$, being the composite of the $C^\infty$ map $t\mapsto e^{t}$
(with values in $(0,1)$) with $y$, which is $C^\infty$ on $(0,1)$.  By the chain
rule, with $r=e^{t}$,
\[
  Y'(t)=e^{t}y'(e^{t})=r\,y'(r),
\]
and differentiating once more,
\[
  Y''(t)=e^{t}y'(e^{t})+e^{2t}y''(e^{t})=r\,y'(r)+r^{2}y''(r).
\]
Since $r\bigl(r\,y'\bigr)'(r)=r\bigl(y'(r)+r\,y''(r)\bigr)=r\,y'(r)+r^{2}y''(r)$,
\begin{equation}\tag{1}
  Y''(t)=r\bigl(r\,y'\bigr)'(r),\qquad Y'(t)=r\,y'(r)\qquad\bigl(r=e^{t}\bigr).
\end{equation}
Every $r\in(0,1)$ equals $e^{t}$ for exactly one $t\in J$, namely $t=\log r$, so
the hypothesis together with (1) says precisely that
\begin{equation}\tag{2}
  Y''(t)+2n\,Y'(t)+n^2\,Y(t)=0\qquad\text{for every }t\in J .
\end{equation}

\emph{Step 2: a function with vanishing derivative on $J$ is constant.}  We
isolate this step, used twice below, because the standard statement is for a
closed interval while $J$ is open.  Let $g\colon J\to\mathbb{C}$ be
differentiable with $g'(t)=0$ for every $t\in J$, and let $t_0,t_1\in J$ with
$t_0\le t_1$.  Since $J$ is an interval, $[t_0,t_1]\subseteq J$; $g$ is
continuous on the compact interval $[t_0,t_1]$, and for every $x\in[t_0,t_1)$
the derivative of $g$ at $x$ within $[x,\infty)$ exists and equals $0$, being a
restriction of the two-sided derivative $g'(x)=0$.  The theorem that a function
continuous on $[t_0,t_1]$ whose derivative within $[x,\infty)$ vanishes at every
$x\in[t_0,t_1)$ is constant on $[t_0,t_1]$ gives $g(t_1)=g(t_0)$.  As any two
points of $J$ are comparable, $g$ is constant on $J$.

\emph{Step 3: first integrating factor.}  Put
\[
  W(t)=Y'(t)+n\,Y(t)\qquad(t\in J),
\]
which is $C^\infty$ on $J$.  By (2),
\[
  W'(t)=Y''(t)+n\,Y'(t)=-2n\,Y'(t)-n^2Y(t)+n\,Y'(t)
  =-n\bigl(Y'(t)+n\,Y(t)\bigr)=-n\,W(t).
\]
Hence $g(t)=e^{nt}W(t)$ is differentiable on $J$ with
\[
  g'(t)=e^{nt}\bigl(W'(t)+n\,W(t)\bigr)=0\qquad(t\in J),
\]
so by Step 2 it is constant; call its value $B\in\mathbb{C}$.  Thus
\begin{equation}\tag{3}
  Y'(t)+n\,Y(t)=W(t)=B\,e^{-nt}\qquad(t\in J).
\end{equation}

\emph{Step 4: second integrating factor --- this is where the logarithm
enters.}  Let
\[
  h(t)=e^{nt}Y(t)-B\,t\qquad(t\in J),
\]
which is differentiable on $J$.  Using (3),
\[
  \frac{d}{dt}\Bigl(e^{nt}Y(t)\Bigr)=e^{nt}\bigl(Y'(t)+n\,Y(t)\bigr)
  =e^{nt}\cdot B\,e^{-nt}=B ,
\]
so $h'(t)=B-B=0$ for every $t\in J$.  By Step 2, $h$ is constant on $J$; call
its value $A\in\mathbb{C}$.  (Note what has happened: because the two roots
coincide, the right-hand side of (3) after multiplication by $e^{nt}$ is the
\emph{constant} $B$ rather than a nonconstant exponential, and its primitive is
the linear function $B\,t$.  This is the single point at which the argument
differs from the non-degenerate one, and it is what produces the logarithm.)
Hence
\begin{equation}\tag{4}
  e^{nt}Y(t)=A+B\,t,\qquad\text{i.e.}\qquad Y(t)=\bigl(A+B\,t\bigr)e^{-nt}
  \qquad(t\in J).
\end{equation}

\emph{Step 5: back to the variable $r$.}  Let $r\in(0,1)$ and put $t=\log r$.
Then $t<0$, so $t\in J$, and $e^{t}=r$.  For a real exponent $\lambda$ and
$r>0$ one has $e^{\lambda\log r}=r^{\lambda}$, so $e^{-nt}=r^{-n}$.  Therefore
by (4)
\[
  y(r)=Y(\log r)=\bigl(A+B\log r\bigr)r^{-n},
\]
which is the assertion, with the constants $A$ and $B$ produced in Steps 4
and 3.
\end{proof}
